\begin{textsample}{POS Dim 3 – pi – Score 14.00 – pi/pi\_inf\_7.txt}  \label{ex:f3_pos_257}
6.
TRANSIÇÃO DA EDUCAÇÃO \textbf{INFANTIL} AO LONGO DA ETAPA A educação \textbf{infantil} integra a Educação Básica sendo a primeira etapa desse processo educacional.
Com isso, o Ensino Fundamental também passou por adequações ampliando -se para nove anos.
Dessa forma, o currículo da educação \textbf{infantil} pode contribuir para uma visão sistêmica da educação integrando -se com os Anos Iniciais do ensino fundamental.
Para isso é preciso tornar efetiva a progressão de uma etapa para outra através da continuidade e transição educativa como articulação curricular.
Para Serra ( 2004 ), a sequência educacional se refere à maneira como estão organizados os saberes, de forma sequenciada e organizada ao longo dos vários níveis e etapas no processo educacional, considerando o desenvolvimento e aprendizagens das \textbf{crianças}.
Essa progressão entre a Educação \textbf{Infantil} e o Ensino Fundamental não se reduz à sobreposição de uma etapa em relação a outra, nem mesmo que a Educação \textbf{Infantil} seja uma preparação para o ensino fundamental.
Nesse sentido, os professores precisam trabalhar de maneira consciente a efetividade desse processo progressivo.
O processo sistemático de ensino induz ao discente uma mudança em sua vida, na capacidade de pensar, refletir e compreender.
Essa relação subjetiva que o aluno adquire conduz a uma capacidade maturacional constituída no \textbf{dia} a \textbf{dia}, mediada pelos \textbf{estímulos}.
A Educação \textbf{Infantil} como etapa inicial, prepara o discente para o processo de ensino, fomentando sua autonomia, interação e socialização, como instrumento de conduzir o aluno ao convívio em espaço social e escolar.
Pensa, nesse primeiro momento, na transição da \textbf{criança} se inserindo em um mundo novo, pessoas novas, que parece ser muito difícil.
O processo de adaptação às novas ações e comportamentos sem sombra de dúvida vai necessitar de um grande esforço dos educadores para a construção do vínculo bem construído e amenizar o sentimento de medo e angústia da \textbf{criança}.
Na primeira fase de transição, muito mais que adaptação, será um momento de inserção da \textbf{criança} na \textbf{instituição}.
Gradativamente ela vai construindo confiança e vínculo com o professor, aprendendo a lidar com a ausência e a presença da família, iniciando a construção de novos vínculos.
A mudança de uma etapa para outra na Educação \textbf{Infantil} sempre acontece mediante a preparação da \textbf{criança} para cada etapa seguinte.
A autonomia da \textbf{criança} vai acontecendo desde quando ela aprende a conhecer e reconhecer de forma contínua os objetos e ações didáticas trabalhadas.
Essa progressão deve ser monitorada pelos professores através de portfólios, \textbf{registros} e observações das aprendizagens das \textbf{crianças} de maneira contínua, pois é preciso um encadeamento, ou seja, uma sequência lógica no conhecimento construído pelas \textbf{crianças}.
Fonte : Elaborado pelo autor, baseado em Freire, 2017.
Para garantir que os objetivos de aprendizagens priorizem as interações e \textbf{brincadeiras}, os professores precisam sempre retomá -los para ter certeza que as experiências que eles propõem às \textbf{crianças} contemplam tais objetivos.
As estratégias para garantir o sucesso dos trabalhos com foco no alcance desses objetivos devem contemplar a \textbf{roda} de conversa, a convivência com outras \textbf{crianças} para \textbf{brincar} e interagir.
As \textbf{crianças} devem escolher as atividades favoritas para \textbf{brincar} livremente e explorar os materiais, dentre outras estratégias que viabilizem o trabalho dos professores.
Na imagem ao lado pode -se perceber que as transições acontecem em toda a educação básica.
As mudanças de uma fase para outra trazem momentos únicos, em especial a preocupação docente que se dá na transição entre essas duas etapas da Educação Básica, conforme a BNCC/2017, ela requer muita atenção para que haja \textbf{equilíbrio} entre as mudanças introduzidas, garantindo integração e continuidade dos processos de aprendizagens das \textbf{crianças}, respeitando suas singularidades e as diferentes relações que elas estabelecem com os conhecimentos, assim como a natureza das mediações de cada etapa.
Torna -se necessário estabelecer estratégias de acolhimento e adaptação tanto para as \textbf{crianças} quanto para os docentes, de modo que a nova etapa se construa com base no que a \textbf{criança} sabe e é capaz de fazer, em uma perspectiva de continuidade de seu percurso educativo.
O currículo do Estado do Piauí reconhece que a Educação \textbf{Infantil} é uma etapa essencial e avança no sentido de que a \textbf{criança} deve estar no centro do processo de aprendizagem.
O referido documento leva os professores a olhar para as particularidades das \textbf{crianças}, nas três faixas \textbf{etárias}, onde deverão se apropriar do conhecimento e de novas experiências.
As atividades nas \textbf{creches} e nas \textbf{instituições} de ensino precisam ganhar uma nova dinâmica objetivando reforçar a relação de \textbf{afeto} e confiança entre professores e \textbf{crianças}.
De acordo com a BNCC o Currículo de Educação \textbf{Infantil} do Estado do Piauí alinha -se aos 6 ( seis ) direitos de aprendizagens e aos 5 ( cinco ) campos de experiências que estão na BNCC.
Dessa forma, o planejamento dos professores precisa estar alinhado a esses direitos e campos de experiências.
Nesse sentido, o currículo apresenta uma nova organização que coloca a \textbf{criança} como centro e protagonista do processo de ensino aprendizagem.
Esse novo currículo traz conceitos importantes como \textbf{cuidar} e educar com foco nas potencialidades e experiências de cada \textbf{criança}.
Com a implementação da BNCC, e com o currículo da EI nas escolas, vem junto as transformações e perspectivas do mundo contemporâneo.
Com isso, a escola também precisa se atualizar, pois a sociedade exige uma nova forma de trabalhar o processo ensino aprendizagem com foco em campos de experiências e direitos de aprendizagem, possibilitando a construção da autonomia dos estudantes.
Além da reprodução, numa escala ampliada das múltiplas habilidades nas quais a atividade produtiva não poderia ser realizada, o complexo sistema educacional da sociedade é também responsável pela produção e reprodução da estrutura de valores dentro da qual os indivíduos definem seus próprios objetivos e fins específicos.
As relações sociais de produção capitalistas não se perpetuam automaticamente ( MÉSZÁROS, 1981, p. 260 ).
A educação não se desenvolve na fragmentação de conteúdos e na memorização.
Atualmente a escola deve pautar seus trabalhos nas transformações do mundo e na contextualização da comunidade escolar.
Entretanto, é imprescindível que haja no contexto da escola, espaços educativos que atendam as demandas através das interações e \textbf{brincadeiras} necessárias a essa etapa de ensino.

% matched lemmas: afeto, brincadeira, brincar, creche, criança, cuidar, dia, equilíbrio, estímulo, etário, infantil, instituição, registro, roda
\end{textsample}
